% r03_fill_K5b: proof of (K_0), part VIII: the skew-torus instances and the counterterm bridge.

\subsubsection*{The skew-torus instances}

Call a torus $\R^3/\Pi$ \emph{aligned} in the frame $(f_\alpha)=(Oe_\alpha)$ of a lattice $O(s\Z^3)$ if $\Pi=O(\ell_1\Z\times\ell_2\Z\times\ell_3\Z)$ for some $\ell_\mu>0$, that is, if the torus is a rectangular box with periodic boundary conditions whose sides are parallel to the frame vectors, and \emph{skew} in that frame otherwise. \cite{Ba1,Ba2,Ba3} and \cite{K1} are written for a torus aligned in the frame of the blocks (\cite[(2.21)]{K1}). $T_1=\DJ$ is aligned (a cube) in the frame of $\Xk$ and skew in the frame of $\Yk$; $T_2=R^{-1}\DJ$ is aligned for $\Yk$ and skew for $\Xk$. On $T_1$ every block-spin step performed in $Y$'s frame, that is all $r$ steps of the standard flow and steps $2,\dots,r$ of the hybrid flow, is a step on a torus skew in the frame of its blocks; on $T_2$ only the first hybrid step is.

\begin{proposition}[the skew-torus instances]\label{rot:prop:SK}
King's Theorem~3.1 \cite[(3.4)]{K1} (Ba{\l}aban's upper bound with the renormalisation of \cite{Ba3}) and Ba{\l}aban's lower-bound machinery \cite[\S\S2--4, (1.14)]{Ba1} hold for the standard flow of $\Yk$ on $T_1$, the torus $R\DJ$ in $Y$'s frame, which is skew, with the constants of the aligned case times at most $125$ in the linear part; and Ba{\l}aban's lower-bound argument holds for steps $2,\dots,r$ of the hybrid flow on $T_1$ with the constant $C(1+k)^p$ of \Cref{rot:cor:OF}. In both, the torus enters \cite{Ba1,Ba2,Ba3} only through the embeddings of the block lattices, the external fields and the image terms of the periodised kernels; no torus symmetry and no torus Fourier analysis is used in the proof of \cite[Prop.~1]{Ba3}. The statements hold under the provisos $J\in5^{n_*}\N$ and $J\ge J_0(\delta_r)$ of \Cref{rot:thm:CT}.
\end{proposition}

The proof consists of three transfer lemmas. Let $T'\defeq R\DJ$ be the skew torus in $Y$'s frame with frame period lattice $\Gamma'\defeq R^{\mathrm T}(50MJ\,\Z^3)$ (read in the frame coordinates), its cubic cover $\tilde T\defeq\R^3/250MJ\,\Z^3$ (a cube aligned in $Y$'s frame, side $250MJ\in M\N$), the covering $\pi:\tilde T\to T'$ with deck group $\cG=\Gamma'/250MJ\,\Z^3$ of order $125$, and the sup-systole $s_\infty(\Gamma')=40MJ$ (Ba{\l}aban's metric is the sup metric, \cite[p.~605]{Ba1}).

\begin{lemma}[folding]\label{rot:lem:F}
$250MJ\,R^{-1}\Z^3=50MJ\,(5R^{-1})\Z^3\subset50MJ\,\Z^3=P_J$ because $5R^{-1}$ is integral, with index $|\det5R^{-1}|=125$; so $\tilde T$ covers $T'$ with a deck group of order $125$. Every block lattice $Y_{k-j}$ contains $P_J$ (\Cref{rot:lem:adm}(iii)), so the deck translations commute with every block average and the Gaussian quadratic forms of the flow are deck-invariant; $\pi$ preserves bonds, blocks, big blocks and the contours, $\dist_{\tilde T}(\tilde x,\pi^{-1}S)=\dist_{T'}(x,S)$ in the quotient sup metric, and lifted operators intertwine: $G_j(\pi^{-1}\Omega,A\circ\pi)(f\circ\pi)=(G_j(\Omega,A)f)\circ\pi$, likewise $D^\eta$, the H\"older quotients, $\delta G_j$, $C^{(j)}$. Hence every kernel of a $Y$-frame flow on $T'$ ($G_j$, $\Delta^{(j)}$, $C^{(j)}$, the slice kernels) is the $125$-fold periodisation of the corresponding kernel on the aligned torus $\tilde T$, and \cite[Props.~2.1--2.3]{Ba1} ($=$\cite{Ba4}) and \cite[Thm.~3.3]{K1} hold on $T'$ with the $\tilde T$ constants times at most $125$ (fibre points are at least $s_\infty$ apart, and the finite sum of exponentially decaying terms is absorbed).
\end{lemma}
\begin{proof}
The covering step of the aligned case with $(\Z/3)^3$ replaced by $\cG$: a non-wrapping parallelepiped lifts to disjoint parallelepipeds \cite[p.~611]{Ba1}, and positivity on $T'$ is positivity on $\cG$-invariant functions. The deck group commutes with the tile average as well ($P_J\subset Y_{k-1}$), so the lemma covers the hybrid steps.
\end{proof}

\begin{lemma}[restriction principle]\label{rot:lem:D}
Call a functional $\Phi$ liftable of degree $a\in\{0,1\}$ if $\Phi_{\tilde T}(\xi\circ\pi)=125^a\Phi_{T'}(\xi)$ (local sums $a=1$, sup-type $a=0$, quadratic forms in Ba{\l}aban's operators on unions of big blocks $a=1$, background fields $a=0$). If an inequality $\Phi\le\Psi$ between liftable functionals of the same degree holds on $\tilde T$ for all configurations obeying block-local restrictions $\chi$ and all admissible regions, it holds on $T'$ with the same constants. Moreover $\ell^2(T')$ is the $\cG$-invariant subspace of $\ell^2(\tilde T)$ (inner product times $125$) and lifted operators commute with $\cG$, so $\|X_{T'}\|\le\|\tilde X\|$ and $\tilde X\ge0\Rightarrow X_{T'}\ge0$. The lemma does not carry integrals over field configurations (the $\cG$-invariant model on $\tilde T$ is a different measure); those are handled by \Cref{rot:lem:TC}.
\end{lemma}
\begin{proof}
Lifted configurations obey the lifted restrictions; lifted regions $\pi^{-1}\Lambda$ are unions of big blocks with the same distance and corridor data (\Cref{rot:lem:F}); quadratic forms scale by $125$ by the intertwining; backgrounds are images of intertwined operators (the model is abelian). Evaluate the $\tilde T$ inequality at $\xi\circ\pi$ and divide by $125^a$.
\end{proof}

\begin{lemma}[type closure]\label{rot:lem:TC}
With $\sigma_j\defeq s_\infty(\Gamma')/(5^j\eps)$, every step of \cite[\S\S2--4]{Ba1}, \cite[\S\S2--3]{Ba2}, \cite[\S\S1--3]{Ba3} and \cite[\S3]{K1} is of one of six types and holds verbatim on $T'$ with the aligned constants: \textup{(I)} exact identities on a finite graph (Gaussian integration, composition of the block-spin maps, block normalisations, set decompositions); \textup{(II)} pointwise bounds on the integrand (dropping non-negative terms, Neumann decoupling, indicator $\le$ exponential, rescalings); \textup{(III)} operator inequalities valid on any finite graph (the diamagnetic inequality, $\log\det(1+X)\le\Tr X$, $\Tr(X\mathbf 1_\Lambda)\le\|X\|\,|\Lambda|$, $-\Delta_A\ge0$, also with Neumann restriction); \textup{(IV)} decay and norm bounds for $G_j$, $C_j$, $\Delta^{(j)}$, $H_j$, $\delta G_j$ on unions of large blocks, wrapping allowed (\Cref{rot:lem:F,rot:lem:D}); \textup{(V)} local arguments on a set of sup-diameter $D$ (gauge-outs, Stokes, local expansions, the contour conventions inside a block, ``sets of sup-diameter $<\sigma/2$ do not wrap''), verbatim when $D<\sigma_j$: the sets are blocks, big blocks and corridors whose width in scale-$j$ units grows only logarithmically in $t_j^{-1}$, and since the product of the spacing with the corridor width increases along the flow the binding case is the coarsest scale $\delta_r$, which is the condition $40MJ\ge\delta_r\max\{6\sqrt3\,r(\delta_r)+2,\,2/\delta_0\}$, i.e.\ $J\ge J_0(\delta_r)$; \textup{(VI)} counts ($|T|$, $|\Lambda|$, $|\partial\Lambda|$, sums over scales, and the estimate \cite[(3.42)]{Ba2} in the form of \Cref{thm:a06:BALA}(d), whose box partition descends to $T'$ under the same non-wrapping condition; the printed cube count \cite[(3.45)]{Ba2} and the first inequality of (3.46) are false as printed, at every block factor, and are not used): they depend on the torus only through the block count. The large-field analysis (\cite[\S3]{Ba2}, \cite[\S4]{Ba1}, \cite[\S3.2]{K1}) uses (II)--(III) on large-field blocks, (I), (III), (IV) for the fluctuation integral with Neumann decoupling, (VI) for the sum over large-field configurations, $\prod_j\sum_\Lambda e^{-p(t_j)|\Lambda|}\le e^{C|T|}$, and (I)/(IV)/(V) with the perturbative estimates for the collar terms; the lower bound \cite[\S4]{Ba1} restricts to the whole-torus small-field region by Jensen's inequality and uses (I)--(IV) only. The torus enters the prints only through admissibility, the systole, the kernels and their images, counts, counterterms and Wick images, external-field periodisation and symmetries, and once through a coordinate-parity split of the bonds inside \cite[Prop.~3.1$'$]{Ba4} ($=$\cite[(3.29)]{Ba2}), a quadratic-form statement that folds from the cover (\Cref{rot:lem:D}; on $T'$ the periods are even at every scale, so it also holds verbatim); the one momentum-space formula of \cite{Ba3}, (3.28) on p.~441, uses the cubic symmetry with an image-term defect $e^{-c\sigma}$.
\end{lemma}
\begin{proof}
Each display of \cite[Ch.~1--4]{Ba1}, \cite[pp.~556--594]{Ba2}, \cite[pp.~411--445]{Ba3} and \cite[App.]{K1}, read from the page images, is assigned to a type or to one of the entry points listed; types (I)--(III) are torus-free, (IV) is \Cref{rot:lem:F,rot:lem:D}, (V) is the non-wrapping condition, and (VI) is the block count. The one external lemma, the cumulant bound, is \Cref{rot:thm:R1} on the skew tori.
\end{proof}

\begin{proof}[Proof of \Cref{rot:prop:SK}]
The linear part is \Cref{rot:lem:F}; the nonlinear part (the small- and large-field analysis) is \Cref{rot:lem:TC} relative to the same bounds on aligned tori at block factor $5$, which are \Cref{thm:a06:BALA}. The hybrid steps $2,\dots,r$ act on composite blocks with the hybrid kernels of \Cref{rot:prop:BD}(a); Ba{\l}aban's Neumann cube step is replaced by the constant-background reduction of \Cref{rot:lem:CB} and the whole-torus form of \cite[Prop.~1]{Ba3} for the hybrid flow is \Cref{rot:cor:OF}, whence the constant $C(1+k)^p$. The provisos: $J\in5^{n_*}\N$ because the frame period lattice $\Gamma'$ must lie on the $5$-adic grid $M_b5^{n_*}\Z^3$, $n_*$ the number of large-field scales of the $Y$ flow, fixed by the couplings; and the non-wrapping cap $J\ge J_0(\delta_r)$ of type (V). Both are harmless and cofinal: in the flow of this subsection the coarsest spacing is $\delta_r\to0$ (\Cref{rot:lem:r}), so no condition on $J$ remains for large $k$.
\end{proof}

\subsubsection*{The counterterm bridge}

\begin{proposition}[the counterterm and Wick-constant conventions]\label{rot:prop:CT}
\cite{K1} and \cite{Ba1,Ba2,Ba3} are proved for the model in which the mass counterterm is defined by \cite[(3.33)]{K1} with the sum over the torus and torus propagators in the graphs, the Wick subtraction by \cite[(2.8)]{K1} with the torus propagator, the vacuum counterterm $E_1$ by \cite[(3.35)]{K1} as vacuum graphs on the torus, and the source vertex by \cite[(3.21)]{K1} with torus propagators. \Cref{rot:def:tori} uses the free-boundary conventions of \cite[p.~330]{K2}. For a torus $T$ with shortest period $\ell(T)$ ($=50MJ$ for $T_1$, $T_2$) and spacing $s$, write $\dm_T(s)$, $E_{1,T}(s)$, $c_T(s)$ for the torus-convention quantities and
\[
\Delta_T(s)\defeq\dm(s)-\dm_T(s),\qquad\Delta E_T(s)\defeq E_1(s)-E_{1,T}(s),\qquad\Delta c_T(s)\defeq\sW_s-c_T(s)=-\sum_{\gamma\in P\setminus0}C_s(\gamma).
\]
\begin{enumerate}[label=\textup{(\alph*)},leftmargin=*]
\item With $N=2$,
$\Zun_{\Lambda,T}(G,h)=e^{-\Delta E_T|T|-N\Delta c_T\sS_\Lambda(h)}\int\prod dA\prod d\phi\ e^{-S^{\rm tor}_\Lambda-\frac12\Delta_T\sum_xs^3|\phi(x)|^2+F_\Lambda(G)+W^{\rm tor}_\Lambda(h)}$:
the free-boundary model is the torus-convention model with the bare mass shifted by $\Delta_T$ and its counterterms frozen at $m^2$, times a $c$-number and a shift of the Wick constant. Both block-spin flows carry these $c$-numbers through unchanged, so the same identity holds between the effective actions at every step.
\item $\Delta_T(s)$, $\Delta c_T(s)$ and $\Delta E_T(s)$ are bounded uniformly in $s\le s_0$, and $|\Delta_T(s)-\Delta_T(0)|\le Cs$. Put $\Delta_*\defeq\sup_{s\le s_0}|\Delta_T(s)|$ and $u(s)\defeq-\Delta_T(s)/2$. The free-boundary effective polynomial of \Cref{rot:def:eff} equals King's torus-convention polynomial at the constant-shifted source $h+u(s)\mathbf 1$, minus a $c$-number, up to graphs that carry at least one $u$-vertex and have order $\ge14$ in the couplings, whose sum is bounded by $C\delta_r^\sigma|T|$ on $\supp\chi_r$; and \Cref{rot:prop:Kforward}, the only consumer of \Cref{rot:prop:H3,rot:prop:H4,rot:prop:H5}, is $c$-number blind: multiplying each lattice's $\Zun$ and small-field integral by an arbitrary positive constant leaves every step and the conclusion unchanged. Consequently the extensive bounds \Cref{rot:prop:H3,rot:prop:H4,rot:prop:H5}, proved for the torus-convention model at the sources $(zG,h+u(s)\mathbf 1)$, hold for the model of \Cref{rot:def:tori} as consumed; for \Cref{rot:prop:H3,rot:prop:H4} the constant-source shift needs King's and Ba{\l}aban's bounds with constants uniform over the family $\{h+u\mathbf 1:|u|\le\Delta_*/2\}$, which is the uniformity of \cite[Thm.~3.1]{K1} in the scalar mass over a compact mass set, King's own clause ``$C$ depends on $g$, $h$ and $|T|$''; for \Cref{rot:prop:H5} the shift is removable by Jensen's inequality.
\item The ratios are convention-free up to a shift of the source by a constant: for every $\Lambda$, $G$, $h$,
\[
\Zun^{\rm free}_{\Lambda,T}(G,h)=\exp\big[-(\Delta E_T(s)+\Delta_T(s)c_T(s))|T|-2\Delta c_T(s)\sS_\Lambda(h)\big]\,\Zun^{\rm tor}_{\Lambda,T}\big(G,h-\tfrac12\Delta_T(s)\mathbf 1\big),
\]
with a prefactor independent of $G$; so the walled curvature ratio $\langle e^{tF(G)}\rangle_{\fc,\Lambda}$ is the same number in both conventions at the shifted wall $h_0-\Delta_T(\eps_k)/2$.
\end{enumerate}
\end{proposition}
\begin{proof}
(a) and (c) are exact algebra: the two actions differ by the counterterms $\frac12\Delta_T\sum_xs^3|\phi|^2+\Delta E_T|T|$ and the Wick squares by the $c$-number $-2\Delta c_T\sS_\Lambda(h)$; a field-independent constant passes through every block-spin step (\Cref{lem:a06:CF}). (b) For $\Delta c_T=-\sum_{\gamma\ne0}C_s(\gamma)$ the bound is the exponential decay of the massive lattice propagator, $|C_s(\gamma)|\le Ce^{-c|\gamma|}$ uniformly in $s$, summed over the non-zero periods. For $\Delta_T$, in each of the finitely many graphs of \cite[(3.33)]{K1} (orders $a+2b\le4$) the torus-minus-free difference replaces one propagator by its image sum, which is smooth and removes the coincident-point singularity, so that the remaining singular factors are summable in three dimensions uniformly in $s$; this gives the boundedness, and the rate $|\Delta_T(s)-\Delta_T(0)|\le Cs$ follows at one loop from the explicit integrals and at the two-loop orders $(0,2)$, $(2,1)$, $(4,0)$ from a graph-by-graph computation of the thirteen two-loop topologies at order $(4,0)$, the degree-zero loops evaluated inner-first and symmetrised, using that the $e^2\lambda$ vertex-type loop vanishes by charge conjugation and the exchange composed with the photon self-energy vanishes by lattice transversality. For $\Delta E_T$ the same image-sum argument applies to the vacuum graphs of \cite[(3.35)]{K1}, with one vertex fixed. The statement on the effective polynomials: a constant shift of the scalar source enters King's expansion as a vertex $u\sum\eta^3|\phi|^2$ with a bounded coefficient, and the graphs with a $u$-vertex and order $\le13$ in the couplings are exactly the terms by which the two polynomials differ beyond the $c$-number, so the remainder is a finite set of graphs of order $\ge14$ with running power $\ge7>3$, bounded by $C\delta_r^\sigma|T|$. The $c$-number blindness of \Cref{rot:prop:Kforward} is read off its proof.
\end{proof}
